\chapter{The primitive seed and the typed operator grammar}\label{ch:seed-first}

\section{Why the seed must come before the coordinates}
The calculus begins before a projector is chosen.  A projector tells us how to resolve an object that already exists; it does not explain how that object was generated, which operations are admissible, or which relations make two operator expressions represent the same value.  The first layer therefore separates six pieces of data:
\[
\boxed{
\Seed=(X,S,\mathcal A,\mathcal R,\mathcal P,\tau).
}
\]
Here $X$ is the ambient type, $S$ is the seed family, $\mathcal A$ is the fixed operator grammar, $\mathcal R$ is the relation system, $\mathcal P$ is a packet atlas, and $\tau$ is the declared completion structure.  In a finite realization, $\tau$ is automatic; in an infinite realization it must be named because algebraic generation and topological completion are different operations.

\SafeFigure[.98\linewidth]{38_seed_to_reconstruction_pipeline.pdf}{The complete dependency path.  The later seven-area examples start at the left and do not skip directly to a convenient coordinate representation.}{fig:seed-full-pipeline}

\begin{boundarybox}[ against hidden encoding]
A datum is not a useful seed merely because a sufficiently complicated operator can decode the target from it.  The grammar must be fixed independently of the target, the relations must be explicit, and the generated hull must be measurable.  Otherwise the phrase ``one seed generates the object'' has no mathematical content.
\end{boundarybox}

\section{Typed operator words}
A typed grammar may be represented by a quiver.  Vertices are object types and arrows are admissible primitive operations.  A composable path
\[
\gamma=a_m\cdots a_2a_1
\]
is an operator word.  A realization assigns a vector space or module $X_v$ to each vertex and a linear map to each arrow.  Evaluation sends a path to the corresponding composition.

For the one-space finite core used below, all arrows act on one space $X$, so the grammar generates a unital algebra
\[
\mathcal A\subseteq\operatorname{End}_K(X).
\]
The point of the typed formulation is that the same definitions remain meaningful when different arrows have different source and target spaces.

\begin{definition}[Generated hull]
For a seed family $S\subseteq X$, define
\[
\Hull_{\mathcal A}(S)
=
\operatorname{span}_K\{As:A\in\mathcal A,\ s\in S\}.
\]
When $\mathcal A$ is given by primitive generators, $A$ ranges over all finite admissible operator words and their finite linear combinations.
\end{definition}

\begin{theorem}[Smallest invariant subspace]\label{thm:smallest-hull}
$\Hull_{\mathcal A}(S)$ is the smallest $\mathcal A$-invariant linear subspace of $X$ containing $S$.
\end{theorem}

\begin{proof}
The hull contains $S$ because the identity belongs to the unital algebra.  If $B\in\mathcal A$ and $x=\sum_jc_jA_js_j$ lies in the hull, then
\[
Bx=\sum_jc_j(BA_j)s_j\in\Hull_{\mathcal A}(S),
\]
so the hull is invariant.  Conversely, any invariant subspace containing $S$ contains every $As$ and therefore their span.
\end{proof}

\SafeFigure[.93\linewidth]{39_seed_system_anatomy.pdf}{The seed system keeps ambient type, primitive data, grammar, relations, packet coordinates, and completion data separate.}{fig:seed-anatomy}

\section{Seed rank}
\begin{definition}[Seed rank]
The seed rank of a finite $\mathcal A$-module $X$ is
\[
\operatorname{srank}_{\mathcal A}(X)
=
\min\{|S|:\Hull_{\mathcal A}(S)=X\}.
\]
A one-seed realization is called cyclic.
\end{definition}

For a single operator $U\in\operatorname{End}(X)$ and one proposed seed $s$, define the orbit matrix
\[
\mathcal O_U(s)=\begin{bmatrix}s&Us&U^2s&\cdots&U^{n-1}s\end{bmatrix},
\qquad n=\dim X.
\]

\begin{theorem}[Finite cyclicity criterion]\label{thm:orbit-rank}
The vector $s$ generates $X$ under the polynomial algebra $K[U]$ if and only if
\[
\rank\mathcal O_U(s)=n.
\]
\end{theorem}

\begin{proof}
The columns of the orbit matrix span the Krylov space generated by $s$.  Cayley--Hamilton reduces every power $U^m$, $m\ge n$, to a linear combination of $I,U,\ldots,U^{n-1}$.  Thus the full generated hull equals the column space of $\mathcal O_U(s)$, and this hull is $X$ precisely when the matrix has rank $n$.
\end{proof}

\section{Master seed: one cyclic shift}
We now introduce the single example that will be carried through every layer of the calculus.  Let
\[
X=\C^8,
\qquad
s=e_0=(1,0,0,0,0,0,0,0)^{\mathsf T},
\]
and let $U$ be the forward cyclic shift
\[
Ue_j=e_{j+1\pmod 8}.
\]
Its matrix has ones immediately below the diagonal and one in the upper-right corner.

\begin{calculationbox}[ orbit]
The first eight operator words give
\[
s=e_0,
\quad Us=e_1,
\quad U^2s=e_2,
\quad\ldots,\quad
U^7s=e_7.
\]
Therefore
\[
\mathcal O_U(s)=I_8
\]
up to the chosen ordering of the standard basis, and
\[
\boxed{\rank\mathcal O_U(s)=8.}
\]
The seed is cyclic and $\Hull_{\C[U]}(s)=X$.
\end{calculationbox}

\SafeFigure[.78\linewidth]{40_seed_orbit_matrix.pdf}{The orbit columns are the eight coordinate vectors.  The rank calculation is exact and does not depend on a numerical tolerance.}{fig:seed-orbit-matrix}

\SafeFigure[.82\linewidth]{41_generated_hull_growth.pdf}{The dimension of the partial hull after admitting $s,Us,\ldots,U^ks$.  At depth seven the full eight-dimensional space has been generated.}{fig:hull-growth}

\begin{stepbox}{Seed-to-finish Step 1: primitive datum}
Start with $s=e_0$.  No signal coefficients, Fourier bins, wavelet coefficients, or packet boundaries have yet been inserted.
\end{stepbox}

\begin{stepbox}{Seed-to-finish Step 2: fixed grammar}
Declare one primitive operator $U$ and permit finite polynomial expressions in $U$.  The target vector introduced later must be generated from this already-declared grammar.
\end{stepbox}

\begin{stepbox}{Seed-to-finish Step 3: generated hull}
The orbit rank proves that the declared seed and grammar generate the entire ambient space.  This is the first exact certificate in the chain.
\end{stepbox}

\chapter{Relations, quotients, and finite normal forms}\label{ch:relations}

\section{The free presentation}
Generation alone produces many operator words that may evaluate to the same element.  Relations are the kernel data that turn the free word space into the actual realized object.

Let $\mathcal A^{(S)}$ denote the free left $\mathcal A$-module on formal symbols $[s]$, $s\in S$.  Define the evaluation map
\[
\varepsilon:\mathcal A^{(S)}\longrightarrow X,
\qquad
\varepsilon\!\left(\sum_{s\in S}A_s[s]\right)
=
\sum_{s\in S}A_ss.
\]

\begin{theorem}[Seed presentation theorem]\label{thm:seed-presentation}
The evaluation map has image $\Hull_{\mathcal A}(S)$ and induces an isomorphism
\[
\boxed{
\mathcal A^{(S)}/\Ker\varepsilon
\cong
\Hull_{\mathcal A}(S).
}
\]
\end{theorem}

\begin{proof}
The image is exactly the span of the values $As$, so it is the generated hull.  The quotient isomorphism is the first isomorphism theorem.
\end{proof}

\section{The shift relation}
For the master example,
\[
U^8=I.
\]
Consequently the polynomial evaluation map
\[
\varepsilon_s:\C[z]\longrightarrow X,
\qquad
p(z)\longmapsto p(U)s
\]
has $(z^8-1)$ in its kernel.

\begin{proposition}[Exact cyclic quotient]\label{prop:cyclic-quotient}
For the cyclic shift seed,
\[
\boxed{
X\cong \C[z]/(z^8-1)
}
\]
as a $\C[z]$-module, where multiplication by $z$ corresponds to $U$.
\end{proposition}

\begin{proof}
The map is surjective by cyclicity.  Both the quotient and $X$ have dimension eight.  Since $(z^8-1)$ lies in the kernel, the induced surjection from the eight-dimensional quotient to $X$ is an isomorphism.
\end{proof}

\SafeFigure[.94\linewidth]{42_relation_quotient.pdf}{The exact relation $U^8=I$ converts infinitely many formal words into eight normal-form coordinates.}{fig:relation-quotient}

\section{Normal form and the first nontrivial object}
Every generated object has a unique representative
\[
p(U)s
=
\sum_{k=0}^7a_kU^ks.
\]
Choose
\[
p(z)=2-z+3z^2+z^4+4z^5-2z^6+5z^7.
\]
Then
\[
\boxed{
x=p(U)s=(2,-1,3,0,1,4,-2,5)^{\mathsf T}.
}
\]
The eight entries have not been hidden in an arbitrary decoder: they are the coefficients of the unique normal form in the fixed quotient algebra.

\begin{calculationbox}[ coefficient extraction]
Because $U^ks=e_k$,
\[
\begin{aligned}
p(U)s
&=2e_0-e_1+3e_2+0e_3+e_4+4e_5-2e_6+5e_7\\
&=(2,-1,3,0,1,4,-2,5)^{\mathsf T}.
\end{aligned}
\]
\end{calculationbox}

\begin{stepbox}{Seed-to-finish Step 4: relation system}
Record $U^8=I$.  This is the precise rule that allows arbitrarily long operator words to be reduced to a finite normal form.
\end{stepbox}

\begin{stepbox}{Seed-to-finish Step 5: generated object}
Choose a polynomial $p$ after the grammar and relation are fixed, and form $x=p(U)s$.  The object now exists before any packet atlas is chosen.
\end{stepbox}

\section{Symbolic certificate}
The exact implementation is short because the mathematical normal form has already been proved.
\begin{lstlisting}
import sympy as sp
from proof_atlas.foundations import cyclic_shift

U = cyclic_shift(8)
s = sp.eye(8)[:, 0]
coeff = [2, -1, 3, 0, 1, 4, -2, 5]
x = sum((coeff[k] * U**k * s for k in range(8)), sp.zeros(8, 1))

assert U**8 == sp.eye(8)
assert sp.Matrix.hstack(*[U**k * s for k in range(8)]).rank() == 8
assert x == sp.Matrix(coeff)
\end{lstlisting}

\chapter{Addresses before projectors}\label{ch:addresses}

\section{Address maps}
An address system is a declared rule assigning each elementary coordinate to a packet label.  Let $\{e_j\}_{j\in J}$ be a basis and let
\[
a:J\longrightarrow A
\]
be an address map.  Define
\[
P_\alpha e_j=
\begin{cases}
e_j,&a(j)=\alpha,\\
0,&a(j)\ne\alpha.
\end{cases}
\]

\begin{theorem}[Address-projector theorem]\label{thm:address-projector}
The operators constructed from an address map satisfy
\[
P_\alpha P_\beta=\delta_{\alpha\beta}P_\alpha,
\qquad
\sum_{\alpha\in A}P_\alpha=I.
\]
Conversely, every complete family of pairwise annihilating projectors gives the direct-sum decomposition
\[
X=\bigoplus_{\alpha\in A}P_\alpha X.
\]
\end{theorem}

\begin{proof}
Each basis vector belongs to exactly one address class, so applying two distinct selectors kills it and summing all selectors returns it.  The identities follow on a basis and hence on all of $X$.  Conversely, completeness gives $x=\sum_\alpha P_\alpha x$, and pairwise annihilation implies uniqueness of this decomposition.
\end{proof}

\section{Three nested address levels}
For the master vector, use three levels:
\[
\begin{array}{ccl}
\text{level 1} &:& j\pmod 2,\\
\text{level 2} &:& j\pmod 4,\\
\text{level 3} &:& j\pmod 8.
\end{array}
\]
Let $E_j$ be the individual coordinate projector.  Let
\[
R_r=E_r+E_{r+4},\qquad 0\le r<4,
\]
and
\[
P_{\mathrm e}=E_0+E_2+E_4+E_6,
\qquad
P_{\mathrm o}=E_1+E_3+E_5+E_7.
\]
Then
\[
\boxed{
P_{\mathrm e}=R_0+R_2,
\qquad
P_{\mathrm o}=R_1+R_3,
\qquad
R_r=E_r+E_{r+4}.
}
\]

\SafeFigure[.98\linewidth]{43_address_refinement_tree.pdf}{The address tree is declared before the numerical packets are formed.  Every parent is exactly the sum of its descendants.}{fig:address-tree}

\section{Root-of-unity addresses from the seed operator}
The same cyclic operator that generated the object also generates its Fourier packet projectors.  Let $\omega=e^{2\pi i/N}$ and suppose $U^N=I$.  Define
\[
Q_k
=
\frac1N\sum_{m=0}^{N-1}\omega^{-km}U^m,
\qquad 0\le k<N.
\]

\begin{theorem}[Cyclic character projectors]\label{thm:character-projectors}
The operators $Q_k$ satisfy
\[
Q_kQ_\ell=\delta_{k\ell}Q_k,
\qquad
\sum_{k=0}^{N-1}Q_k=I,
\qquad
UQ_k=\omega^kQ_k
\]
with the sign convention determined by the chosen Fourier transform.
\end{theorem}

\begin{proof}
Multiply the two sums and collect the coefficient of $U^r$.  The inner character sum is
\[
\sum_{m=0}^{N-1}\omega^{(\ell-k)m}
=
N\delta_{k\ell}.
\]
This proves idempotence and annihilation.  Summing $Q_k$ over $k$ leaves only the $m=0$ term, giving $I$.  Multiplication by $U$ shifts the summation index and produces the character factor.
\end{proof}

\begin{remark}
The spatial packets $E_j$ and Fourier packets $Q_k$ are not competing truths.  They are two exact atlases of the same generated object.  Their usefulness is determined by the operator or question being studied.
\end{remark}

\begin{stepbox}{Seed-to-finish Step 6: addresses}
Declare parity, residue-mod-four, coordinate, and Fourier-character addresses.  Different atlases may coexist because each has its own exact synthesis law.
\end{stepbox}

\SafeFigure[.96\linewidth]{44_analysis_synthesis_packets.pdf}{The master vector resolved into its eight coordinate packets and reassembled.  Negative packet values remain exact coordinates; no positivity assumption is needed.}{fig:master-spatial-packets}
